\documentclass[12pt,a4paper]{article}
\usepackage[utf8]{inputenc}
\usepackage[T1]{fontenc}
\usepackage[spanish]{babel}
\usepackage{geometry}
\usepackage{hyperref}
\usepackage{booktabs}
\usepackage{xcolor}
\usepackage{enumitem}
\usepackage{tikz}
\usetikzlibrary{arrows.meta}

\geometry{margin=2.5cm}

\title{\textbf{Reporte de Avance del Pipeline EPM: DLC + YOLO BBox Constraint + SimBA}}
\author{TT\_Ratones\_2026}
\date{19 de marzo de 2026}

\begin{document}

\maketitle

\section{Resumen Ejecutivo}
Durante esta iteraci\'on se corrigi\'o el principal cuello de botella del sistema: la calidad de la pose generada por DeepLabCut antes de entrar a SimBA. El problema hist\'orico eran los keypoints err\'aticos, especialmente cuando se analizaban videos con \textit{downscale}. Esos errores aplanaban o descompon\'ian las predicciones de \textit{Grooming} y \textit{Thigmotaxis}.

La mejora principal fue sustituir el suavizado gen\'erico por un post-proceso guiado por YOLO, llamado \textbf{YOLO BBox Constraint}. Adem\'as:
\begin{itemize}[leftmargin=1.5cm]
    \item se elimin\'o el \textit{downscale} del flujo nuevo,
    \item se dejaron s\'olo \textbf{6 paredes} como ROIs can\'onicas de SimBA para \textit{Thigmotaxis},
    \item se regener\'o el set supervisado con pose filtrada,
    \item se reentren\'o \textit{Grooming},
    \item y se entren\'o de forma correcta \textit{Thigmotaxis} en modo \texttt{Train single model (global environment)}.
\end{itemize}

El resultado cualitativo pas\'o de ``malo'' a ``decente'' en videos nuevos. Todav\'ia hay falsos positivos y omisiones, especialmente en \textit{Grooming}, pero la calidad general del sistema subi\'o de forma clara y reproducible.

\section{Problema Detectado}
Los problemas originales eran:
\begin{itemize}[leftmargin=1.5cm]
    \item keypoints de DLC que se teletransportaban fuera del cuerpo del rat\'on,
    \item degradaci\'on visible cuando se aplicaba \textit{downsampling},
    \item features de SimBA contaminadas por una pose anat\'omicamente imposible,
    \item respuestas pobres o planas en los indicadores finales de \textit{Grooming} y \textit{Thigmotaxis}.
\end{itemize}

\section{Arquitectura Actual del Pipeline}
El flujo operativo actualizado es:

\begin{center}
\texttt{Video original o clip recortado}
$\rightarrow$
\texttt{DeepLabCut full-res}
$\rightarrow$
\texttt{YOLO BBox Constraint}
$\rightarrow$
\texttt{pose corregida}
$\rightarrow$
\texttt{compute\_simba\_features.py}
$\rightarrow$
\texttt{SimBA}
$\rightarrow$
\texttt{video multimodal final}
\end{center}

\subsection{Diagrama del Pipeline}

\begin{center}
\begin{tikzpicture}[
    >=Latex,
    box/.style={
        draw,
        rounded corners,
        align=center,
        minimum width=3.2cm,
        minimum height=1.1cm,
        font=\small,
        fill=blue!6
    },
    decision/.style={
        draw,
        rounded corners,
        align=center,
        minimum width=3.4cm,
        minimum height=1.1cm,
        font=\small\bfseries,
        fill=green!8
    },
    simba/.style={
        draw,
        rounded corners,
        align=center,
        minimum width=3.4cm,
        minimum height=1.1cm,
        font=\small,
        fill=orange!10
    },
    out/.style={
        draw,
        rounded corners,
        align=center,
        minimum width=3.4cm,
        minimum height=1.1cm,
        font=\small,
        fill=red!8
    }
]
\node[box] (video) at (0,0) {Video original\\o clip recortado};
\node[box] (dlc) at (4.4,0) {DeepLabCut\\full-resolution};
\node[decision] (bbox) at (8.9,0) {YOLO BBox\\Constraint};
\node[box] (pose) at (13.4,0) {Pose corregida\\H5 / CSV / MP4};

\node[simba] (bridge) at (4.4,-2.4) {compute\_simba\_features.py\\bridge 8bp + input\_csv};
\node[simba] (simba) at (8.9,-2.4) {SimBA\\features + clasificadores};
\node[out] (hud) at (13.4,-2.4) {Video multimodal\\trayectoria + m\'etricas};

\draw[->, thick] (video) -- (dlc);
\draw[->, thick] (dlc) -- (bbox);
\draw[->, thick] (bbox) -- (pose);
\draw[->, thick] (pose) |- (bridge);
\draw[->, thick] (bridge) -- (simba);
\draw[->, thick] (simba) -- (hud);
\end{tikzpicture}
\end{center}

\subsection{Sin Downscale}
El pipeline nuevo ya no usa \textit{downscale} cuando el video se analiza desde cero. Se mantiene el recorte temporal, pero el clip de trabajo conserva la resoluci\'on original para DLC.

\subsection{Filtro YOLO BBox Constraint}
No se est\'a usando Savitzky-Golay como filtro principal. Se implement\'o un filtro personalizado basado en YOLO:
\begin{enumerate}[leftmargin=1.5cm]
    \item YOLO detecta al rat\'on por frame y genera un \textit{bounding box}.
    \item Cada keypoint de DLC se valida contra el \textit{bbox} con margen.
    \item Los puntos imposibles se invalidan.
    \item La pose se interpola.
    \item Se hace una segunda pasada espacial con:
    \begin{itemize}[leftmargin=1.5cm]
        \item revalidaci\'on contra bbox,
        \item l\'imite de salto por keypoint,
        \item radio m\'aximo respecto al centro corporal,
        \item y \textit{snap} final al \'ultimo punto v\'alido o al centro estimado.
    \end{itemize}
\end{enumerate}

Esto no s\'olo suaviza la trayectoria: impone una restricci\'on anat\'omica basada en la detecci\'on visual del animal.

\subsection{Validaci\'on Visual}
El pipeline genera un MP4 de inspecci\'on con:
\begin{itemize}[leftmargin=1.5cm]
    \item el video original,
    \item la pose corregida,
    \item el \textit{bounding box} de YOLO,
    \item y el contador de keypoints corregidos.
\end{itemize}

Esto permite que el investigador valide la calidad de la pose antes de confiar en SimBA.

\section{Integraci\'on de Zonas y ROIs}
La configuraci\'on de zonas desde Streamlit qued\'o separada en dos capas:

\subsection{Diagrama de Capas Espaciales}

\begin{center}
\begin{tikzpicture}[
    >=Latex,
    layer/.style={
        draw,
        rounded corners,
        minimum width=4.8cm,
        minimum height=1.2cm,
        align=center,
        font=\small
    },
    note/.style={
        draw,
        rounded corners,
        minimum width=4.2cm,
        minimum height=0.9cm,
        align=center,
        font=\small
    }
]
\node[layer, fill=purple!10] (streamlit) at (0,0) {Streamlit\\dibujo de zonas y paredes};
\node[layer, fill=cyan!10] (yolo) at (-4.2,-2.6) {YOLO / visualizaci\'on\\Brazo abierto, brazo cerrado, centro};
\node[layer, fill=orange!12] (simba) at (4.2,-2.6) {SimBA / thigmotaxis\\Pared 1 a Pared 6};
\node[note, fill=green!10] (result) at (0,-5.0) {Resultado:\\menos saturaci\'on de ROIs y mejores features espaciales};

\draw[->, thick] (streamlit) -- (yolo);
\draw[->, thick] (streamlit) -- (simba);
\draw[->, thick] (yolo) -- (result);
\draw[->, thick] (simba) -- (result);
\node[font=\scriptsize] at (-2.5,-1.1) {capa visual};
\node[font=\scriptsize] at (2.5,-1.1) {capa can\'onica};
\end{tikzpicture}
\end{center}

\subsection{ROIs para SimBA}
Para \textit{Thigmotaxis}, SimBA usa s\'olo \textbf{6 paredes}:
\begin{center}
\texttt{Pared 1, Pared 2, Pared 3, Pared 4, Pared 5, Pared 6}
\end{center}

Estas paredes se guardan correctamente en el proyecto SimBA y aparecen como \texttt{ROIs defined}. Son la base espacial can\'onica para las features de contacto/cercan\'ia a pared.

\subsection{Zonas para conteo y visualizaci\'on}
Por separado, en Streamlit se mantienen:
\begin{itemize}[leftmargin=1.5cm]
    \item \texttt{Brazo Abierto 1}
    \item \texttt{Brazo Abierto 2}
    \item \texttt{Brazo Cerrado 1}
    \item \texttt{Brazo Cerrado 2}
    \item \texttt{Centro}
\end{itemize}

Estas zonas son m\'as \'utiles para el conteo por brazos y para el render multimodal que para SimBA.

\section{Rebuild del Set Supervisado}
Se implement\'o un batch para refrescar el dataset etiquetado:
\begin{itemize}[leftmargin=1.5cm]
    \item aplica \texttt{BBox Constraint} a la pose original,
    \item regenera \texttt{features\_extracted},
    \item reconstruye \texttt{targets\_inserted},
    \item conserva las etiquetas supervisadas hist\'oricas de \textit{Grooming} y \textit{Thigmotaxis}.
\end{itemize}

Esto permiti\'o rehacer el set de entrenamiento sobre una pose mucho m\'as limpia.

\section{Entrenamiento de Modelos}
\subsection{Grooming}
Se confirm\'o un nuevo modelo:
\begin{center}
\texttt{models/generated\_models/Grooming.sav}
\end{center}

\subsection{Thigmotaxis}
Se detect\'o primero que el entrenamiento anterior no hab\'ia dejado un \texttt{.sav} nuevo usable. Despu\'es se ejecut\'o correctamente el modo \texttt{Train single model (global environment)} y se confirm\'o:
\begin{center}
\texttt{models/generated\_models/Thigmotaxis.sav}
\end{center}

con fecha del 19 de marzo de 2026.

\subsection{Consumo de Modelos en la App}
Se actualiz\'o la aplicaci\'on para que \texttt{pages/04\_Analisis\_Final.py} y \texttt{src/scripts/generar\_video\_prediccion.py} prefieran:
\begin{itemize}[leftmargin=1.5cm]
    \item \texttt{models/generated\_models/Grooming.sav}
    \item \texttt{models/generated\_models/Thigmotaxis.sav}
\end{itemize}

y usen los modelos viejos de \texttt{validations} s\'olo como respaldo.

\section{Resultados Cualitativos}
El sistema mejor\'o de forma visible:
\begin{itemize}[leftmargin=1.5cm]
    \item las predicciones dejaron de ser planas o completamente muertas,
    \item \textit{Thigmotaxis} responde mejor en videos nuevos con paredes bien configuradas,
    \item \textit{Grooming} ya da se\~nal, aunque sigue siendo la parte m\'as fr\'agil.
\end{itemize}

La conclusi\'on cualitativa actual es:
\begin{center}
\textbf{antes: malo / ahora: decente}
\end{center}

Todav\'ia hay:
\begin{itemize}[leftmargin=1.5cm]
    \item falsos positivos,
    \item grooming no detectado en algunos bouts reales,
    \item grooming activado en momentos biol\'ogicamente dudosos,
    \item y generalizaci\'on a\'un imperfecta fuera del set supervisado original.
\end{itemize}

\section{Lectura T\'ecnica del Estado Actual}
El proyecto ya no est\'a bloqueado por un error estructural. El pipeline funciona y generaliza mejor que antes, pero todav\'ia necesita m\'as variedad de entrenamiento.

La interpretaci\'on m\'as razonable es:
\begin{itemize}[leftmargin=1.5cm]
    \item el mayor problema hist\'orico s\'i era la pose,
    \item el filtro bbox corrigi\'o una parte muy grande del problema,
    \item ahora los errores restantes son m\'as propios del dataset y del reentrenamiento que de la ingenier\'ia del pipeline.
\end{itemize}

\section{Pr\'oximos Pasos}
\begin{enumerate}[leftmargin=1.5cm]
    \item Etiquetar de forma supervisada videos nuevos, empezando por \texttt{R5C\_01mar24}.
    \item Incorporar m\'as videos limpios del nuevo lote entregado por el doctor.
    \item Seguir validando la pose filtrada con MP4s de inspecci\'on.
    \item Mostrar de forma m\'as clara los artefactos finales dentro de \texttt{Resultados y Estad\'isticas}.
    \item Formalizar un esquema de carpetas por corrida final para trazabilidad.
\end{enumerate}

\subsection{Diagrama de Mejora Continua}

\begin{center}
\begin{tikzpicture}[
    >=Latex,
    ring/.style={
        draw,
        rounded corners,
        minimum width=3.3cm,
        minimum height=1.0cm,
        align=center,
        font=\small,
        fill=blue!7
    }
]
\node[ring] (a) at (0,0) {Videos nuevos\\limpios};
\node[ring] (b) at (5.2,0) {DLC + BBox Constraint};
\node[ring] (c) at (5.2,-2.3) {Anotaci\'on\\supervisada};
\node[ring] (d) at (0,-2.3) {Reentrenamiento\\SimBA};

\draw[->, thick] (a) -- (b);
\draw[->, thick] (b) -- (c);
\draw[->, thick] (c) -- (d);
\draw[->, thick] (d) -- (a);
\end{tikzpicture}
\end{center}

\section{Conclusi\'on}
Este avance es importante porque transform\'o el proyecto de un pipeline inestable y enga\~noso a un pipeline funcional y auditable. Todav\'ia falta mejorar la calidad del aprendizaje supervisado y ampliar el dataset, pero la base t\'ecnica correcta ya est\'a puesta:
\begin{itemize}[leftmargin=1.5cm]
    \item pose full-resolution,
    \item filtro anat\'omico guiado por YOLO,
    \item ROIs correctas para \textit{Thigmotaxis},
    \item modelos nuevos en \texttt{generated\_models},
    \item y mejor respuesta en videos nuevos.
\end{itemize}

\end{document}
